%==============================================================================
\section{Histograms --- Biểu đồ tần suất}
\label{sec:histograms}
%==============================================================================

\begin{dinhnghia}[title={Histogram}]
Histogram là biểu diễn dạng cột (giống biểu đồ cột) của \textbf{phân phối} một biến.
Nó cho biết \textbf{tần suất} (số lần xuất hiện) của từng giá trị (hoặc khoảng giá trị).
Trục $x$: miền giá trị; trục $y$: tần suất hoặc số đếm; \textbf{chiều cao cột} = số điểm rơi vào khoảng đó.
\end{dinhnghia}

\begin{ynghia}[title={Histogram dùng để làm gì?}]
Histogram hữu ích khi nhận diện pattern: \textbf{độ lệch} (skewness), \textbf{số đỉnh} (modality), \textbf{điểm kỳ dị} (outlier).
\begin{itemize}
  \item \textbf{Skewness}: mức bất đối xứng của phân phối.
  \item \textbf{Modality}: số ``đỉnh'' trên phân phối.
  \item \textbf{Outlier}: điểm nằm ngoài vùng giá trị điển hình.
\end{itemize}
\end{ynghia}

\subsection{Histogram trong Python}

\begin{phuongphap}[title={Thư viện và dataset ví dụ}]
Python có Matplotlib, Seaborn, Plotly, Bokeh.
Ví dụ dưới dùng \textbf{Matplotlib} và dataset \textbf{Breast Cancer} từ scikit-learn: 30 feature, 569 mẫu; mô tả đặc tính tế bào ung thư vú và nhãn ác tính/lành tính.
\end{phuongphap}

\subsection{Ví dụ --- Histogram cơ bản}

\begin{vidu}[title={Import và histogram mean radius}]
\begin{lstlisting}
import matplotlib.pyplot as plt
from sklearn.datasets import load_breast_cancer
data = load_breast_cancer()
\end{lstlisting}

\begin{lstlisting}
import matplotlib.pyplot as plt
from sklearn.datasets import load_breast_cancer
data = load_breast_cancer()

plt.figure(figsize=(7.2, 3.5))
plt.hist(data.data[:,0], bins=20)
plt.xlabel('Mean Radius')
plt.ylabel('Frequency')
plt.show()
\end{lstlisting}

Dùng \texttt{hist()} của Matplotlib; \texttt{bins=20} chia miền giá trị thành 20 khoảng; \texttt{xlabel}/\texttt{ylabel} gắn nhãn trục.
\end{vidu}

\begin{ynghia}[title={Đọc biểu đồ}]
Phân phối mean radius gần \textbf{chuẩn}, đỉnh khoảng 12--14.
\end{ynghia}

\begin{vidu}[title={Output}]
\tsimg{02_histograms/img01.jpg}
\end{vidu}

\subsection{Ví dụ --- Histogram nhiều tập dữ liệu}

\begin{ynghia}[title={So sánh hai nhóm}]
Có thể chồng hai histogram để so sánh phân phối --- ví dụ mean radius giữa mẫu ác tính và lành tính.
\end{ynghia}

\begin{vidu}[title={Hai histogram chồng nhau}]
\begin{lstlisting}
import matplotlib.pyplot as plt
from sklearn.datasets import load_breast_cancer
data = load_breast_cancer()

plt.figure(figsize=(7.2, 3.5))
plt.hist(data.data[data.target==0,0], bins=20, alpha=0.5, label='Malignant')
plt.hist(data.data[data.target==1,0], bins=20, alpha=0.5, label='Benign')
plt.xlabel('Mean Radius')
plt.ylabel('Frequency')
plt.legend()
plt.show()
\end{lstlisting}

Gọi \texttt{hist()} hai lần; \texttt{alpha=0.5} để cột trong suốt, không che khuất hoàn toàn; \texttt{legend()} thêm chú thích.
Trong scikit-learn, \texttt{target==0} là ác tính (malignant), \texttt{target==1} là lành tính (benign) --- theo nhãn dataset.
\end{vidu}

\begin{ynghia}[title={Nhận xét}]
Hai phân phối \textbf{khác nhau}: nhóm ác tính có tần suất cao hơn ở mean radius lớn.
\end{ynghia}

\begin{vidu}[title={Output}]
\tsimg{02_histograms/img02.jpeg}
\end{vidu}
